\begin{ouverture}[calcul]{3}{Multiples, diviseurs et partages}{Partie 1 · Nombres entiers et calculs}
{Critères · Listes de diviseurs · Communs}
\objectif{Reconnaître multiples et diviseurs, avec les critères par 2, 5 et 10.}
\objectif{Trouver tous les diviseurs d’un nombre jusqu’à 30.}
\objectif{Chercher des diviseurs ou des multiples communs.}
\prerequis{Tables de multiplication et division euclidienne.}
\end{ouverture}
\pagestyle{otMain}\essentiel\fontsize{11}{15}\selectfont
\phantomsection\label{cours-3-1}
\essSection{1}{Multiples et diviseurs}
\begin{definition}Comme $56=8\times7$, le nombre 56 est un multiple de 8 et de 7 ; et 8 et 7 sont des diviseurs de 56 : le partage tombe juste, le reste de la division est nul.\end{definition}
\begin{propriete}Critères sur le chiffre des unités : multiple de 2 s’il finit par 0, 2, 4, 6, 8 ; multiple de 5 s’il finit par 0 ou 5 ; multiple de 10 s’il finit par 0. Pour les autres diviseurs, on utilise les tables ou une division.\end{propriete}
\begin{exemple}7 divise-t-il 91 ? $7\times13=91$ : oui. 6 divise-t-il 92 ? $92=6\times15+2$ : non, le reste vaut 2.\end{exemple}
\begin{savoirfaire}[Lire le chiffre des unités]Un nombre proche peut aider : 91 est à côté de 90, multiple de 5 et de 10 ; mais 91 finit par 1 : il n’est multiple ni de 2, ni de 5, ni de 10. Les critères se lisent sur le nombre lui-même.\end{savoirfaire}
\clearpage
\phantomsection\label{cours-3-2}
\essSection{2}{Trouver tous les diviseurs}
\begin{savoirfaire}[Des paires de facteurs]Pour lister tous les diviseurs de 24, chercher les paires de facteurs : $1\times24$, $2\times12$, $3\times8$, $4\times6$. La liste complète est 1, 2, 3, 4, 6, 8, 12, 24 : rien n’est oublié, car les paires épuisent les cas.\end{savoirfaire}
\begin{exemple}Diviseurs de 18 : 1, 2, 3, 6, 9, 18. Diviseurs de 24 : 1, 2, 3, 4, 6, 8, 12, 24. Les diviseurs communs aux deux nombres sont 1, 2, 3 et 6 : ils permettent des partages identiques des deux collections.\end{exemple}
\begin{propriete}Les multiples communs servent aux rendez-vous réguliers : les multiples de 6 sont 6, 12, 18, 24… ; ceux de 8 sont 8, 16, 24… Le plus petit multiple commun non nul de 6 et 8 est 24.\end{propriete}
\begin{attention}1 divise tous les nombres, et chaque nombre est son propre plus grand diviseur. Une liste de diviseurs se vérifie : chaque candidat doit donner un reste nul.\end{attention}

\clearpage\rubrique{Découverte}
\begin{blocLivre}{ex-3-1}\begin{exercice}[1]{Multiples à repérer}
Écrire les cinq premiers multiples non nuls de 7, puis dire si 84 est un multiple de 7.
\end{exercice}\end{blocLivre}
\begin{blocLivre}{ex-3-2}\begin{exercice}[1]{Les critères des unités}
Parmi 230, 345, 1\,308 et 95, lesquels sont multiples de 2 ? de 5 ? de 10 ?
\end{exercice}\end{blocLivre}
\begin{blocLivre}{ex-3-3}\begin{exercice}[1]{Diviseur ou pas ?}
8 est-il un diviseur de 96 ? 7 est-il un diviseur de 93 ?
\end{exercice}\end{blocLivre}
\begin{blocLivre}{ex-3-4}\begin{exercice}[1]{Tous les diviseurs}
Écrire la liste complète des diviseurs de 24, en t’aidant des paires de facteurs.
\end{exercice}\end{blocLivre}
\begin{blocLivre}{ex-3-5}\begin{exercice}[1]{Pairs et impairs}
Sans calculer, dire si chaque somme est paire ou impaire : $346+58$ ; $127+41$ ; $230+87$.
\end{exercice}\end{blocLivre}

\rubrique{Entraînement}
\begin{blocLivre}{ex-3-6}\begin{exercice}[2]{Diviseurs de 30 et 18}
Lister les diviseurs de 30, puis ceux de 18, puis leurs diviseurs communs.
\end{exercice}\end{blocLivre}
\begin{blocLivre}{ex-3-7}\begin{exercice}[2]{Multiples communs}
Écrire les multiples de 6 jusqu’à 50, puis ceux de 8. Quel est leur plus petit multiple commun non nul ?
\end{exercice}\end{blocLivre}
\begin{blocLivre}{ex-3-8}\begin{exercice}[2]{Un chiffre à choisir}
Trouver tous les chiffres $\square$ pour que $47\square$ soit un multiple de 5. Et pour qu’il soit un multiple de 10 ?
\end{exercice}\end{blocLivre}
\begin{blocLivre}{ex-3-9}\begin{exercice}[2]{Le reste raconte}
La division de 125 par 6 donne $125=6\times20+5$. 6 est-il un diviseur de 125 ? Quel est le plus proche multiple de 6 au-dessus de 125 ?
\end{exercice}\end{blocLivre}
\begin{blocLivre}{ex-3-10}\begin{exercice}[2]{Vrai ou faux ?}
« Tout multiple de 10 est aussi un multiple de 5. » « Tout multiple de 5 est aussi un multiple de 10. » Justifier chaque réponse.
\end{exercice}\end{blocLivre}

\rubrique{Problèmes}
\begin{blocLivre}{ex-3-11}\begin{exercice}[3]{Les bouquets}
Une fleuriste a 24 roses et 18 tulipes. Elle veut des bouquets tous identiques, sans fleur restante. Combien de bouquets peut-elle faire ? Donner toutes les possibilités, puis la composition avec 6 bouquets.
\end{exercice}\end{blocLivre}
\begin{blocLivre}{ex-3-12}\begin{exercice}[3]{Les deux phares}
Un phare s’allume toutes les 12 secondes, un autre toutes les 9 secondes. Ils viennent de s’allumer ensemble. Dans combien de secondes s’allumeront-ils de nouveau ensemble ?
\end{exercice}\end{blocLivre}
\begin{blocLivre}{ex-3-13}\begin{exercice}[3]{Les rangées du verger}
Un jardinier plante 90 arbres en rangées égales. Peut-il faire des rangées de 5 arbres ? de 6 ? de 7 ? Donner à chaque fois le nombre de rangées ou le reste.
\end{exercice}\end{blocLivre}
\begin{blocLivre}{ex-3-14}\begin{exercice}[3]{Le nombre secret}
Je suis un diviseur de 24, je suis pair, je suis un multiple de 3 et je suis plus petit que 10. Qui suis-je ? Expliquer pourquoi la réponse est unique.
\end{exercice}\end{blocLivre}
\begin{blocLivre}{ex-3-15}\begin{exercice}[3]{Toujours vrai ?}
Milo affirme : « La somme de trois multiples de 5 est toujours un multiple de 5. » Tester sur un exemple, puis expliquer avec des paquets de cinq.
\end{exercice}\end{blocLivre}

\clearpage
